\chapter{Estrategia de pruebas}\label{cap:pruebas}

\section{Pirámide de pruebas}

Se siguen tres niveles:

\begin{enumerate}[leftmargin=*,itemsep=0.2em]
  \item \textbf{Unitarias.} Cubren funciones puras y módulos
        independientes (agregación FEVER, validadores, utilidades).
  \item \textbf{Integración.} En el backend, prueban endpoints REST
        con la base de datos sustituida por un mock o por una
        \emph{fixture} \emph{in-memory}.
  \item \textbf{Componente / UI.} En el frontend, React Testing
        Library renderiza páginas completas con servicios mockeados.
\end{enumerate}

\section{Backend (Pytest)}

\begin{itemize}[leftmargin=*,itemsep=0.2em]
  \item \path{fakenews-backend/tests/test_smoke.py}: \texttt{/health},
        endpoints públicos básicos.
  \item \path{tests/test_billing.py}: lógica de cuotas y planes.
  \item Tests del módulo \path{fever/}: extracción, recuperación
        (con \emph{stub}), inferencia, agregación, agente \emph{end-to-end}
        contra mocks.
\end{itemize}

Se ejecutan con \texttt{pytest -q} y configuración en
\path{pytest.ini}.

\section{Frontend (Vitest + RTL)}

\begin{itemize}[leftmargin=*,itemsep=0.2em]
  \item \path{src/pages/__tests__/}: tests de páginas
        (login, register, dashboard, admin).
  \item \path{src/lib/__tests__/}: utilidades (i18n, validadores,
        \emph{access control}).
  \item Mocks de servicios para aislar UI/estado.
\end{itemize}

Se ejecutan con \texttt{npm run test}.

\section{Linting y formato}

\begin{itemize}[leftmargin=*,itemsep=0.2em]
  \item Frontend: ESLint con configuración en
        \path{eslint.config.js}, integrado con Vite.
  \item Backend: \texttt{ruff} para \emph{lint}, \texttt{black}
        opcional para formato.
\end{itemize}

\section{Cobertura}

Pytest con \texttt{pytest --cov} reporta cobertura por módulo.
Se establece como objetivo $\geq$ 70\,\% en módulos críticos
(\texttt{billing}, \texttt{fever}). El detalle por archivo se
incluye en el anexo.
